\begin{frame}{Antecedentes}
\begin{itemize}
    \item De 2019 a 2025 el gasto primario aumentó 1,2 puntos del PIB, mientras los ingresos totales disminuyeron 2,4 puntos del PIB frente a 2019.
    \item A pesar de múltiples reformas tributarias, los ingresos tributarios fueron 14,7\% del PIB en 2025, después del repunte transitorio de 2023.
\end{itemize}
\note{[Sources] MFMP 2026, capítulo 1, tablas 1.3 y 1.5.}
\end{frame}

\begin{frame}{Antecedentes}
\begin{itemize}
    \item A pesar de la activación de la cláusula de escape, los ingresos siguen sobreestimándose.
    \item El deterioro del balance y el crecimiento de la deuda serán más acelerados de lo oficialmente pronosticado.
    \item A medida que transferencias como el FOME y el FEPC se han venido marchitando, otros rubros, como salud y pensiones, han aumentado.
\end{itemize}
\end{frame}